%%%%%%%%%%%%%%%%%%%%%%%%%%%%%%%%%%%%%%%%%%%%%%%%%%%%%%%%%%%%%%%%%%%%%%%%%%%%%%%%%%
\begin{frame}[fragile]\frametitle{}
\begin{center}
{\Large Introduction to Yoga and Yogic Practices}
\end{center}
\end{frame}

%%%%%%%%%%%%%%%%%%%%%%%%%%%%%%%%%%%%%%%%%%%%%%%%%%%%%%%%%%%
\begin{frame}[fragile]\frametitle{Syllabus}

\begin{itemize}
\item 1.1  Yoga : Etymology, definitions, aim, objectives and misconceptions. 
\item 1.2  Yoga : Its origin, history and development. 
\item 1.3  Guiding principles to be followed by Yoga practitioners.  
\item 1.4  Principles of Yoga (Triguna, Antahkarana-chatustaya, Tri-Sharira/ Panchakosha). 
\item 1.5  Introduction to major schools of Yoga (Jnana, Bhakti, Karma, Patanjali, Hatha). 
\item 1.6  Introduction to Yoga practices for health and well being. 
\item 1.7  Introduction to Shatkarma: meaning, purpose and their significance in Yoga Sadhana. 
\item 1.8  Introduction to Yogic  Sukshma Vyayama,  Sthula Vyayama and Surya Namaskara.  
\item 1.9  Introduction to Yogasana: meaning, principles, and their health benefits. 
\item 1.10  Introduction to Pranayama and Dhyana and their health benefits. 
\end{itemize}
	  
\end{frame}


%%%%%%%%%%%%%%%%%%%%%%%%%%%%%%%%%%%%%%%%%%%%%%%%%%%%%%%%%%%%%%%%%%%%%%%%%%%%%%%%%%
\begin{frame}[fragile]\frametitle{}
\begin{center}
{\Large 1.1 Yoga : Etymology, definitions, aim, objectives and misconceptions}
\end{center}
\end{frame}

%%%%%%%%%%%%%%%%%%%%%%%%%%%%%%%%%%%%%%%%%%%%%%%%%%%%%%%%%%%
\begin{frame}[fragile]\frametitle{Etymology of Yoga}

  \begin{itemize}
    \item Etymology: the study of where words come from (शब्द व्युत्पत्ती शास्त्र).
	\item \textit{योग} (Yoga) comes from the root \textit{युज्} (Yuj): to join or to unite (युज संयोग).
	\item Devotional and Vedantic reading: the individual self (जीवात्मा) joins the universal self (परमात्मा).
    \item Simpler reading: harmony of body, mind and spirit through mastery of body, senses and mind (not one of Panini's senses).
	\item Panini's three senses of युज (Yuj): union (युजिर योगे), concentration (युज समाधौ) and control (युज संयमने).
	\item Patanjali uses the word mainly in the sense of concentration.
	\item The word means both the goal (साध्य) and the practice that leads there (साधना).
  \end{itemize}

\end{frame}

%%%%%%%%%%%%%%%%%%%%%%%%%%%%%%%%%%%%%%%%%%%%%%%%%%%%%%%%%%%
\begin{frame}[fragile]\frametitle{Definitions of Yoga}
      \begin{itemize}
        \item \textbf{पतञ्जलि  योगसूत्र  Patanjali Yoga Sutras}:      योगश्चित्तवृत्तिनिरोधः || १:२ || Yoga is the stilling of the movements of the mind.
        \item \textbf{Yoga Vasishtha योगवशिष्ठ }:          मनः प्रशमनोपायः योग इत्यभिधीयते | Yoga is the skilful means of calming the mind.
        \item \textbf{भगवद्गीता  Bhagavad Gita}:         
		      \begin{itemize}
				\item  योगस्थः कुरु कर्माणि सङ्गं त्यक्त्वा धनञ्जय |   सिद्ध्यसिद्ध्योः समो भूत्वा \textbf{समत्वं योग उच्यते }|| २:४८ || Act while steady in yoga, free of attachment, even in success and failure. Evenness of mind is yoga.
				\item बुद्धियुक्तो जहातीह उभे सुकृतदुष्कृते |      तस्माद्योगाय युज्यस्व \textbf{योगः कर्मसु कौशलम् }|| २:५० || One joined to wisdom sets aside both good and bad deeds here; so take up yoga. Yoga is skill in action.
				\item युक्ताहारविहारस्य युक्तचेष्टस्य कर्मसु |युक्तस्वप्नावबोधस्य \textbf{योगो भवति दु:खहा }||६ .१ ७ || For one who is moderate in food, recreation, effort, sleep and waking, yoga destroys sorrow.
				\end{itemize}
		\item \textbf{Kathopanishad कठोपनिषद २/३/११}: तां योगामिती मन्यतें स्थिरमिन्द्रिय धारणं Holding the senses back, firm and steady, is called yoga.
      \end{itemize}
\end{frame}

% %%%%%%%%%%%%%%%%%%%%%%%%%%%%%%%%%%%%%%%%%%%%%%%%%%%%%%%%%%%
% \begin{frame}[fragile]\frametitle{Definitions of Yoga}
      % \begin{center}
        % \includegraphics[width=\linewidth,keepaspectratio]{ycb1_definitions}

		% {\tiny (Ref: Param Yoga)}		
        % \end{center}

% \end{frame}

% %%%%%%%%%%%%%%%%%%%%%%%%%%%%%%%%%%%%%%%%%%%%%%%%%%%%%%%%%%%
% \begin{frame}[fragile]\frametitle{Meaning and concepts of Yoga}
% Four Ashrams of Vedic Life: Based on karma कर्म and dharma धर्म, Vedic life was divided into four ashrams: 

          % \begin{itemize}
            % \item \textbf{The First Ashrama - brahmacharya } - "ब्रह्मचर्य" or the Student Stage
            % \item \textbf{The Second Ashrama - gruhastha } - "गृहस्थ" or the Householder Stage
            % \item \textbf{The Third Ashrama - vaanaprastha } - "वानप्रस्थ" or the Hermit Stage
            % \item \textbf{The Fourth Ashrama - sanyas } - "संन्यास" or the Wandering Ascetic Stage
          % \end{itemize}
		  
% Purusharthas ( पुरुषार्थ )literally means ``aims of man'' or ``what men live for''. It is the prescription for the right way of living, recognizing various inherent urges and satisfying them for a meaningful life.
          
          % \begin{itemize}
            % \item \textbf{dharma धर्म}: Righteousness
            % \item \textbf{aarth अर्थ}: Wealth or material means
            % \item \textbf{kaam काम}: Satisfaction of sensual desires
            % \item \textbf{moksh मोक्ष}: Freedom or spiritual liberation
          % \end{itemize}
% \end{frame}

% %%%%%%%%%%%%%%%%%%%%%%%%%%%%%%%%%%%%%%%%%%%%%%%%%%%%%%%%%%%
% \begin{frame}[fragile]\frametitle{The Gurus and Gurukul System}
      % \begin{itemize}
        % \item {\textit “To light the candle one needs another burning candle; in the same way, those who are unenlightened need the help of an illumined Guru.”} - Swami Sivananda
          
          % \item The word Guru:  गु meaning "Darkness" and रु meaning "to destroy". 
		  % \item A Guru is a spiritually enlightened being that dispels the darkness of ignorance and lifts the veil of माया (illusion), rendering the disciple free from the cycle of death and birth.
          
          
          % % The journey stops with the finding of the Guru. The Guru is the gateway to the infinite, guiding the disciple towards the personal realization of the eternal truth.
          
          
          % % The love between Guru and disciple transcends all human relationships. It is through this relationship that the disciple's potential is given shape and direction, working towards the divine plan.
          
        % \item \textbf{The Gurukul System in India}: Gurukul गुरुकुल translates to "the home (family) of the preceptor." 
		% \item The Guru गुरु lives with his family and disciples. Students come to be part of the Guru's family, learning Vedas, Sanskrit, fine arts, administration, and other skills and etiquette.
          
          
          % % Gurukuls were in forests or by riversides. The first 25 years of life, as ब्रह्मचर्य (a celibate student), were spent acquiring training for serving society and personal emancipation.
          
          
          % % Life in Gurukul involved rigorous तपस (austerities) to prepare for life's hardships, instilling habits like hard work, enthusiasm, and selfless service. Students performed duties like fetching water, cleaning, and gathering twigs, while practicing meditation and Yoga.
          
          
          % % Classes were held under shady banyan trees or in thatched huts. Education was generally free, but the Guru received दक्षिणा (fee) from disciples.
          
      % \end{itemize}
% \end{frame}


%%%%%%%%%%%%%%%%%%%%%%%%%%%%%%%%%%%%%%%%%%%%%%%%%%%%%%%%%%%
\begin{frame}[fragile]\frametitle{Aims/Objectives of Yoga}

      \begin{itemize}
		\item Five supports of a healthy, happy life: आहार  Aahar (food), व्यवहार  Vyavahar (behaviour), आचार  Achar (conduct), विचार  Vichar (thought), विहार  Vihar (recreation).
		% the class objectives below are teaching summaries, not a Board list (notes)
		\item To cultivate \textit{Discipline} and \textit{Self-Control}.
		\item To improve \textit{Mental Focus} and \textit{Concentration}.
		\item To enhance \textit{Emotional Stability} and \textit{Resilience}.
		\item To promote \textit{Physical Fitness} and \textit{Posture}.
		\item To achieve \textit{Holistic Well-Being} and \textit{Harmonious Living}.
		% \item लक्ष्य : स्व चे आकलन. आत्मा ते परमात्मा प्रवास 
		% \item उद्देश : सर्वांगीण विकास, सामंजस्य स्थापना. मन बुद्धी व चरित्र यांना शुद्ध बनवणे		
		\item Ultimate aim: freedom from suffering, with a calm mind, will power and a balanced personality (traditionally \textit{kaivalya} or \textit{moksha}).
		\item In public programmes the stated aim is health and well-being.
		% teacher's wording, no named source (notes, ycb1-th-etymology-definitions-aim):
		\item उद्देश्य: सर्वांगीण विकास, सामंजस्य की स्थापना। मन, बुद्धि और चरित्र को शुद्ध बनाना।
	  \end{itemize}

\end{frame}

%%%%%%%%%%%%%%%%%%%%%%%%%%%%%%%%%%%%%%%%%%%%%%%%%%%%%%%%%%%
\begin{frame}[fragile]\frametitle{Misconceptions about Yoga}

      \begin{itemize}
		\item ``Only postures (\textit{aasan आसन })'': no, yoga is holistic, not just exercise.
		\item ``A religion'': no, it belongs to no single religion or community.
		\item ``Needs flexibility first'': no, practice builds it; breath practices need none.
		\item ``Just relaxation'': no, it integrates body, mind and spirit.
		\item ``Only for the young'': no, beginners of all ages can start.
		\item ``Replaces medical treatment'': no, it does not replace medical care.
		% \item  धर्म : फक्त हिंदूंसाठी नाही. वैश्विक. 
		% \item  व्यायाम: फक्त शारीरिक नाही तर मानसिक आणि आध्यत्मिक . 
		% \item  चमत्कार/प्रदर्शन/सिद्धी प्राप्ती 
		% \item  तरुणांसाठीच नाही  तर सर्वांसाठी 
		\item  धर्म: केवल हिंदुओं के लिए नहीं, बल्कि सार्वभौमिक।
		\item  व्यायाम: केवल शारीरिक नहीं, बल्कि मानसिक और आध्यात्मिक भी।
		\item  चमत्कार/प्रदर्शन/सिद्धि प्राप्ति।
		\item  युवाओं के लिए ही नहीं, बल्कि सभी के लिए।		
	  \end{itemize}

\end{frame}


%%%%%%%%%%%%%%%%%%%%%%%%%%%%%%%%%%%%%%%%%%%%%%%%%%%%%%%%%%%%%%%%%%%%%%%%%%%%%%%%%%
\begin{frame}[fragile]\frametitle{}
\begin{center}
{\Large 1.2 Yoga : Its origin, history and development}
\end{center}
\end{frame}


%%%%%%%%%%%%%%%%%%%%%%%%%%%%%%%%%%%%%%%%%%%%%%%%%%%%%%%%%%%
\begin{frame}[fragile]\frametitle{Origin of Yoga}

      \begin{itemize}
		\item Arose in ancient India; the Guidebook says almost 5,000 years ago; exact dates are uncertain.
		\item The Guidebook names the \textit{ऋग्वेद  Rigveda} as the first reference.
		\item One of the षटदर्शन  six Indian philosophical systems (दर्शन ), with \textit{सांख्य  Samkhya} and \textit{वेदान्त  Vedanta}.
        % \item \textbf{Ancient Roots}: 
          
          % Yoga originated in ancient India. The origin of Yoga is traced to the वेद (Vedas). It is difficult to ascertain a fixed time period for the origin of these ancient texts due to differing historical accounts.
          
        % \item \textbf{Ancient Tradition}: 
          
          % The system of Yoga is an ancient tradition having its origin in India. The practice of Yoga is believed to have started at the very dawn of civilization.
          
        \item \textbf{Yogic Lore}: 
          
          In yogic lore, Shiva is the first yogi, आदियोगी (Adiyogi), and the first teacher, आदि गुरु (Adi Guru); he passed it to the seven sages.
          		
		% \item उत्पत्ती: हजारो वर्षांपूर्वी, भारतात. शिव हे आदी योगी आणि गुरु. सप्तर्षींद्वारा प्रसार. 
		% \item प्राचीन पुरातत्व अवशेषांवरून सिद्ध होते कि सिंधू हडप्पा संस्कृतीत योग होता, मुद्रा मुर्त्या होत्या. 
		% \item साहित्य: “हिरण्यगर्भ योगस्य वक्ता मान्य:पुरातन:”, वेद , उपनिषद , दर्शन, बौद्ध, जैन परंपरा 
		% \item पतंजली योगसूत्र 		
		\item उत्पत्ति: हजारों वर्ष पहले, भारत में। शिव आदि योगी और गुरु। सप्तर्षियों द्वारा प्रसार।
		% says the Indus remains "prove" yoga; notes: they are read as signs and historians debate this
		% \item प्राचीन पुरातत्व अवशेषों से सिद्ध होता है कि सिंधु हड़प्पा संस्कृति में योग था, मुद्राएं और मूर्तियां थीं।
		\item Indus Saraswati seals and figures are read as signs of yoga; historians debate this reading.
		% the Hiranyagarbha verse is not in the checked sources (notes, ycb1-th-origin-history-development)
		% \item साहित्य: “हिरण्यगर्भ योगस्य वक्ता मान्य: पुरातन:”, वेद, उपनिषद, दर्शन, बौद्ध, जैन परंपरा।
		\item Literature: Vedas, Upanishads, the darshanas, Buddhist and Jain traditions, then the Yoga Sutra.
		\item  पतंजलि योगसूत्र।
	  \end{itemize}

\end{frame}

%%%%%%%%%%%%%%%%%%%%%%%%%%%%%%%%%%%%%%%%%%%%%%%%%%%%%%%%%%%
\begin{frame}[fragile]\frametitle{Historical Development of Yoga}

      \begin{itemize}
		\item \textbf{Pre-Vedic}: Indus Saraswati seals and figures (about 2700 BC, per the Guidebook) are read as signs of yoga.
		\item \textbf{Vedic}: Vedas and Upanishads; breath, austerity and meditation; ritual turns inward to self-knowledge.
		\item Dates are uncertain; the four period names are the Guidebook's grouping, not fixed years.
		\item \textbf{Classical}: पतञ्जलि (Patanjali) sets out the eightfold path in the \textit{Yoga Sutra}; the \textit{Bhagavad Gita} also belongs here.
		\item \textbf{Post-Classical}: Tantra, then \textit{Hatha Yoga}, drawing on tantric and ascetic traditions; body-centred cleansing, postures, breath control.
		\item \textbf{Modern}: worldwide spread from Swami Vivekananda (Chicago, 1893); then Krishnamacharya, Sivananda, Iyengar, Kuvalayananda and others.

		% removed: invented period dates and unsourced attributions (see reports/ycb1_tools/deck_fixes_intro_faqs.md)
		%
		%
		%
		%
		
	  \end{itemize}

\end{frame}

%%%%%%%%%%%%%%%%%%%%%%%%%%%%%%%%%%%%%%%%%%%%%%%%%%%%%%%%%%%
\begin{frame}[fragile]\frametitle{History of Yoga}
      \begin{center}
        % FIG ycb1-fig-history | now: images/ycb1_history.png | final: books/YCB1/images/new/ycb1_history_new.pdf | register: books/YCB1/docs/image_register.csv
        % PROMPT SP-DIAGRAM (common text: books/YCB1/docs/image_system_prompt.md) | Used in block ycb1-th-origin-history-development, frame 'History of Yoga'. Concept diagram: History of Yoga. Boxes, circles and arrows only, drawn from the structure given in the notes block (the names and their order are listed there); Sanskrit terms in Roman letters with the English meaning in brackets, nothing else added. Flat colours that still read in grayscale. Best redrawn later in TikZ. Aspect ratio 4:3. Check before use: author and the Sanskrit check in the notes.
        % NEW-PIC approved by the author (2026-11-27); the original is commented out below it
        \includegraphics[width=0.7\linewidth,keepaspectratio]{new/ycb1_history_new.pdf}
        % original, replaced (author approved 2026-11-27): \includegraphics[width=0.7\linewidth,keepaspectratio]{ycb1_history}

        % FIG ycb1-fig-teachers | now: images/ycb1_teachers.jpg | final: no image file (see register) | register: books/YCB1/docs/image_register.csv
        % PROMPT SP-NONE (common text: books/YCB1/docs/image_system_prompt.md) | Do not generate portraits of real people. Replace with a plain timeline diagram of the teachers named in the checked notes block, or with photographs the author has permission to use. Used in block ycb1-th-origin-history-development, frame 'History of Yoga'.
        \includegraphics[width=0.7\linewidth,keepaspectratio]{ycb1_teachers}

		{\tiny (Ref: Param Yoga)}		
        \end{center}

\end{frame}

%%%%%%%%%%%%%%%%%%%%%%%%%%%%%%%%%%%%%%%%%%%%%%%%%%%%%%%%%%%
\begin{frame}[fragile,shrink=6]\frametitle{Key Texts and Influences}

  \begin{itemize}
	\item \textbf{Vaidik (वैदिक)}: the early ritual and philosophical base.
	\item \textbf{Upanishad (उपनिषद)}: the concepts of yoga, with stress on self-realisation.
	\item \textbf{Patanjali Yogasutra (पतंजलि योगसूत्र)}: the eightfold path (अष्टांग योग) set out as a system.
	\item \textbf{Bhagavad Gita (भगवद्गीता)}: joins yoga to life and duty (कर्मयोग, भक्तियोग, ज्ञानयोग).
	\item \textbf{Hatha Yoga Pradipika (हठयोग प्रदीपिका)}: the practical methods of hatha yoga, by Svatmarama.
	\item \textbf{भारतीय दर्शन}:
		\begin{itemize}
			\item \textbf{आस्तिक (वेद मानने वाले)}: न्याय (गौतम), वैशेषिक (कणाद), सांख्य (कपिल), मीमांसा (जैमिनी), योग (पतंजलि), वेदांत (बादरायण)।
			\item \textbf{नास्तिक}: जैन (महावीर), बौद्ध (गौतम बुद्ध), चार्वाक (बृहस्पति)।
	    \end{itemize}
	\item \textbf{Samkhya Philosophy}: two eternal principles, Purusha (pure consciousness) and Prakriti (matter).
	\item \textbf{Gunas}: Prakriti is made of three gunas: सत्त्व (purity), रजस (activity) and तमस (inertia).
	\item \textbf{Kaivalya (कैवल्य)}: liberation by telling Purusha apart from Prakriti.
	\item It ends three kinds of suffering: from oneself (adhyatmika), from other beings (adhibhautika), from cosmic forces (adhidaivika).
  \end{itemize}

\end{frame}

%%%%%%%%%%%%%%%%%%%%%%%%%%%%%%%%%%%%%%%%%%%%%%%%%%%%%%%%%%%
\begin{frame}[fragile]\frametitle{Evolution of Yoga Practices}

      \begin{itemize}
		\item Early practice: meditation and austerity.
		\item Patanjali set it out as eight limbs.
		\item Post-classical hatha yoga brought body-centred practice: postures (\textit{आसन}), cleansing, breath control (\textit{प्राणायाम}).
		% the nadi and posture sequence is a summary, not a dated claim (notes)
		\item Traditional schools include \textit{हठ योग}, \textit{कुण्डलिनि योग}, \textit{राज योग} and \textit{कर्म योग}.
		\item Modern named styles: \textit{विन्यास}, \textit{अष्टांग}, Iyengar and \textit{Power Yoga}.
		\item \textbf{सूक्ष्म व्यायाम}: gentle loosening before practice; designed by Maharshi Karthikeyaji Maharaj, spread by स्वामी धीरेंद्र ब्रह्मचारी.
		\item \textbf{स्थूल व्यायाम}: whole-body exercises for energy and strength.

		% \item  सूक्ष्म व्यायाम: हलके, योगाभ्यासाच्या आधी. स्वामी धीरेंद्र ब्रह्मचारी (गुरु: कार्तिकेय महाराज) . शरीर लवचिक व तयार होते. 
		% \item स्थूल व्यायाम: सर्व शरीराचा, गतीचा, शक्तीचा संचार . 
		\item सूक्ष्म व्यायाम: हल्के, योगाभ्यास से पहले। स्वामी धीरेंद्र ब्रह्मचारी (गुरु: कार्तिकेय महाराज)। शरीर लचीला और तैयार होता है।
		\item स्थूल व्यायाम: पूरे शरीर का, गति और शक्ति का संचार।
	  \end{itemize}

\end{frame}



%%%%%%%%%%%%%%%%%%%%%%%%%%%%%%%%%%%%%%%%%%%%%%%%%%%%%%%%%%%%%%%%%%%%%%%%%%%%%%%%%%
\begin{frame}[fragile]\frametitle{}
\begin{center}
{\Large 1.3 Guiding principles to be followed by Yoga practitioners}
\end{center}
\end{frame}

%%%%%%%%%%%%%%%%%%%%%%%%%%%%%%%%%%%%%%%%%%%%%%%%%%%%%%%%%%%%%%%%%%%%%%%%%%%%%%%%%%
\begin{frame}[fragile]\frametitle{Guiding Principles for Yoga Teachers and Practitioners}

	\begin{itemize} 
	\item \textbf{Preparation:} a clean, calm, well-ventilated space; for an online class, check the connection.
	\item \textbf{Diet:} practise on an empty or light stomach; a full stomach makes practice uncomfortable.
	\item \textbf{Starting:} begin with a prayer or invocation.
	\item \textbf{Posture Correction:} ask permission before touching a student.
	\item \textbf{Verbal Cues:} short, clear cues on breath and alignment.
	\item \textbf{Post-Practice:} wait 20 to 30 minutes before a bath or food (Common Yoga Protocol).
	\item \textbf{Feedback:} leave time at the end for feedback and questions.
	\item Online check, touch, cues and feedback are sound teaching practice, not Board rules. \end{itemize}

\end{frame}

%%%%%%%%%%%%%%%%%%%%%%%%%%%%%%%%%%%%%%%%%%%%%%%%%%%%%%%%%%%
\begin{frame}[fragile]\frametitle{Guiding Principles for Yoga Practitioners}

Five points credited to Swami Vishnudevananda, drawn from Swami Sivananda's teaching

    \begin{itemize}
        \item \textbf{Proper Exercise - aasan आसन}
            
            % आसन (Asana) develops the body, broadens mental faculties, and spiritual capacities. Unlike violent physical exercises, yoga postures focus on maintaining the spine’s flexibility and strength.
            
        \item \textbf{Proper Breathing - pranayam प्राणायाम}
            
            % प्राणायाम (Pranayama) involves controlling प्राण (Prana) to achieve control over the mind. Proper breathing increases lung capacity and vitality.
            
        \item \textbf{Proper Relaxation - shavasan  शवासन }
            
            % सवासन (Savasana) is the corpse pose, a state of physical, mental, and spiritual relaxation that recharges and balances the body and mind.
            
        \item \textbf{Proper Diet - Vegetarian}
            
            % A Yogic diet consists of pure, simple, natural foods that promote health. Vegetarian meals aid digestion and provide balanced nutrition.
            
        \item \textbf{Positive Thinking and Meditation - vedant  वेदान्त  and dhyan ध्यान }
            
            % ध्यान (Dhyana) is the practice of meditation that focuses the mind internally on the self, leading to lasting happiness and peace.
            
    \end{itemize}
\end{frame}

%%%%%%%%%%%%%%%%%%%%%%%%%%%%%%%%%%%%%%%%%%%%%%%%%%%%%%%%%%
\begin{frame}[fragile]\frametitle{Practice: Before, During and After}

    \begin{itemize}
        \item \textbf{Before:} keep the place, body and mind clean; calm, quiet place; empty or light stomach; empty bladder and bowels; mat or folded blanket; light, comfortable cotton clothes.
        \item \textbf{During:} begin with a prayer; move slowly, aware of body and breath; do not hold the breath unless the practice says so; breathe through the nose; never jerk or hold the body tight; work within your own capacity; be regular and patient; end with meditation, silence or a peace chant.
        \item \textbf{After:} wait about 20 to 30 minutes before a bath or a meal.
    \end{itemize}

\end{frame}

%%%%%%%%%%%%%%%%%%%%%%%%%%%%%%%%%%%%%%%%%%%%%%%%%%%%%%%%%%
\begin{frame}[fragile]\frametitle{When Not to Practise, and Food}

    \begin{itemize}
        \item Do not practise when exhausted, ill, in a hurry or under acute stress.
        \item Long-term illness, pain or a heart problem: ask a physician or yoga therapist first.
        \item Pregnancy or menstruation: ask a yoga expert first.
        \item Food: a light, mostly vegetarian diet is the usual advice; treat it as advice, not a rule to impose on students.
        \item A little honey in warm water is allowed before practice if a person is weak.
    \end{itemize}

\end{frame}


%%%%%%%%%%%%%%%%%%%%%%%%%%%%%%%%%%%%%%%%%%%%%%%%%%%%%%%%%%%%%%%%%%%%%%%%%%%%%%%%%%
\begin{frame}[fragile]\frametitle{}
\begin{center}
{\Large 1.4 Principles of Yoga (त्रिगुण Triguna, अन्त:करण -चतुस्तय Antahkarana-chatustaya,त्रिशरिर  TriSharira/  पञ्चकोष Panchakosha)}
\end{center}
\end{frame}

%%%%%%%%%%%%%%%%%%%%%%%%%%%%%%%%%%%%%%%%%%%%%%%%%%%%%%%%%%%%%%%%%%%%%%%%%%%%%%%%%%
\begin{frame}[fragile]\frametitle{}
\begin{center}
{\Large Background Terms}
\end{center}
\end{frame}

%%%%%%%%%%%%%%%%%%%%%%%%%%%%%%%%%%%%%%%%%%%%%%%%%%%%%%%%%%%
\begin{frame}[fragile,shrink=6]\frametitle{Understanding Key Concepts}
      \begin{itemize}
        % brahman, maaya and upaadhi are standard Advaita glossary terms, not defined in the Board sources (notes, ycb1-th-principles-of-yoga)
        \item \textbf{brahman ब्रह्मन् }: the absolute reality: eternal, infinite, uncaused; being, consciousness and bliss (सत्-चित्-आनन्द).
        \item \textbf{maaya माया}: the creative power of ब्रह्मन् through which the world appears.
        \item \textbf{atman आत्मन्}: the unchanging inner Self, one with ब्रह्मन्.
        \item \textbf{upaadhi उपाधि}: a limiting adjunct: taking body and mind to be the Self.
        \item \textbf{avidya अविद्या}: ignorance of one's true nature; the first of the five kleshas.
        \item \textbf{jeev जीव}: the individual self as it appears, bound until it knows the आत्मन्.
        \item \textbf{ishwar ईश्वर}: in Patanjali's Sutra, a special purusha, not a religious figure.
        \item \textbf{purush पुरुष}: pure consciousness, in साङ्ख्य-योग.
        \item \textbf{prakruti प्रकृति}: matter, the material cause, made of the three गुण.
      \end{itemize}

\textbf{Summary}: through अविद्या the जीव takes the उपाधि to be itself and suffers, until it knows itself as आत्मन्.
\end{frame}

%%%%%%%%%%%%%%%%%%%%%%%%%%%%%%%%%%%%%%%%%%%%%%%%%%%%%%%%%%%
\begin{frame}[fragile]\frametitle{त्रिगुण  (Three गुण )}

      \begin{itemize}
		\item \textit{sattva सत्व } - Quality of purity, clarity, and harmony.
		\item \textit{rajas रजस } - Quality of activity, movement, and restlessness.
		\item \textit{tamas तमस } - Quality of inertia, darkness, and ignorance.
		\item Balance of गुण  affects mental and emotional states.
		\item Goal: Cultivate \textit{सत्व} for spiritual growth and peace.
	  \end{itemize}

\end{frame}

%%%%%%%%%%%%%%%%%%%%%%%%%%%%%%%%%%%%%%%%%%%%%%%%%%%%%%%%%%%
\begin{frame}[fragile]\frametitle{अन्त:करण -चतुस्तय  (Four Aspects of the Inner Instrument)}

      \begin{itemize}
		\item \textit{manas मनस } (Mind) - Handles thoughts and sensory perceptions.
		\item \textit{buddhi बुद्धि } (Intellect) - Functions as the decision-making faculty.
		\item \textit{ahamkaar अहंकार } (Ego) - Sense of individuality and self-identity.
		\item \textit{chitta चित्त } (Memory) - Stores past experiences and impressions.
		\item Harmonizing these aspects aids in mental clarity and self-awareness.
	  \end{itemize}

\end{frame}

%%%%%%%%%%%%%%%%%%%%%%%%%%%%%%%%%%%%%%%%%%%%%%%%%%%%%%%%%%%
\begin{frame}[fragile]\frametitle{त्रिशरिर  Tri-Sharira (Three Bodies)}

      \begin{itemize}
		\item \textit{sthula sharir स्थूल  शरीर } (Gross Body) - Physical body, the annamaya sheath.
		\item \textit{sukshma sharir सूक्ष्म  शरीर } (Subtle Body) - Made of the pranamaya, manomaya and vijnanamaya sheaths (vital energy, mind, intellect, ego).
		\item \textit{kaaran sharir कारण  शरीर } (Causal Body) - The anandamaya sheath.
		\item Understanding these bodies aids in holistic self-realization.
		\item Goal: Achieve harmony among all three bodies for spiritual growth.
	  \end{itemize}

\end{frame}

%%%%%%%%%%%%%%%%%%%%%%%%%%%%%%%%%%%%%%%%%%%%%%%%%%%%%%%%%%%
\begin{frame}[fragile]\frametitle{पञ्चकोष   (Five Sheaths)}

      \begin{itemize}
		\item \textit{अन्नमय  annamay कोष  kosh} (Food Sheath) - The outermost, physical sheath, made of food.
		\item \textit{प्राणमय  praanamay  कोष  kosh} (Vital Air Sheath) - The energy sheath, a bridge between body and mind.
		\item \textit{मनोमय  manomay कोष kosh } (Mental Sheath) - Mind: thinking and emotion.
		\item \textit{विज्ञानमय Vijnanamaya कोष  kosh} (Wisdom Sheath) - Higher intellect and intuition.
		\item \textit{आनन्दमय Anandamaya कोष  kosh} (Bliss Sheath) - Closest to the true self but still a sheath; the Self lies behind all five.
		\item Practice works outside in: asana on the food sheath, pranayama on energy, meditation on mind and understanding.
	  \end{itemize}

      \begin{center}
        % FIG ycb1-fig-kosha | now: images/ycb1_kosha.png | final: books/YCB1/images/new/ycb1_kosha_new.pdf | register: books/YCB1/docs/image_register.csv
        % PROMPT SP-DIAGRAM (common text: books/YCB1/docs/image_system_prompt.md) | Used in block ycb1-th-principles-of-yoga, frame '(Five Sheaths)'. Concept diagram: (Five Sheaths). Boxes, circles and arrows only, drawn from the structure given in the notes block (the names and their order are listed there); Sanskrit terms in Roman letters with the English meaning in brackets, nothing else added. Flat colours that still read in grayscale. Best redrawn later in TikZ. Aspect ratio 4:3. Check before use: author and the Sanskrit check in the notes.
        % NEW-PIC approved by the author (2026-11-27); the original is commented out below it
        \includegraphics[width=\linewidth,height=0.33\textheight,keepaspectratio]{new/ycb1_kosha_new.pdf}
        % original, replaced (author approved 2026-11-27): \includegraphics[width=0.4\linewidth,keepaspectratio]{ycb1_kosha}

		% credit of the replaced original: {\tiny (Ref: Param Yoga)}
		{\tiny (Line diagram drawn for this deck.)}
        \end{center}
\end{frame}

%%%%%%%%%%%%%%%%%%%%%%%%%%%%%%%%%%%%%%%%%%%%%%%%%%%%%%%%%%%
\begin{frame}[fragile]\frametitle{पञ्च  महाभूत  Pancha Maha Bhuta: Five Elements}


    \begin{itemize}
        \item \textbf{Akasha आकाश  (Space/Ether)}
        \item \textbf{Vayu वायु  (Air)}
        \item \textbf{Agni अग्नि (Fire)}
        \item \textbf{Jala जल (Water)}
        \item \textbf{Prithvi पृथ्वी (Earth)}
    \end{itemize}
	
    \textbf{Elemental Sequence (Taittiriya Upanishad):}
	
    \begin{itemize}
        \item From आकाश (Akasha) comes वायु (Vayu)
        \item From वायु (Vayu) comes अग्नि (Agni)
        \item From अग्नि (Agni) comes जल (Jala)
        \item From जल (Jala) comes पृथ्वी (Prithvi)
    \end{itemize}
	
\end{frame}

%%%%%%%%%%%%%%%%%%%%%%%%%%%%%%%%%%%%%%%%%%%%%%%%%%%%%%%%%%%
\begin{frame}[fragile]\frametitle{पञ्च  महाभूत   Pancha Maha Bhuta: Five Elements}

    \textbf{Physical Properties and Functions:}
	
	
        \begin{tabular}{|p{0.21\linewidth}|p{0.1\linewidth}|p{0.1\linewidth}|p{0.1\linewidth}|p{0.1\linewidth}|p{0.1\linewidth}|}
% \begin{tabular}{|l|l|l|l|l|l|}
    \hline
    \textbf{Element} & \textbf{Space/Ether} & \textbf{Air} & \textbf{Fire} & \textbf{Water} & \textbf{Earth} \\
    \hline
    \textbf{Attributes} & Light, Minute, Quick, Clear, Soft & Light, Dry, Rough, Mobile, Cold, Minute, Clear & Dry, Hot, Sharp, Light, Rough, Minute, Clear & Heavy, Moist, Fluid, Slimy, Cold, Thick, Soft & Heavy, Rough, Solid, Massive, Firm, Hard \\
    \hline
    \textbf{Role in the Body} & Body cavities & Movement, pulsation, conduction & Whole body metabolism & All adhesions, joints & Body organs, mass, inertia \\
    \hline
    \end{tabular}

    {\tiny Property lists differ between sources; this table is not yet checked against the Guidebook's own lists.}

\end{frame}


%%%%%%%%%%%%%%%%%%%%%%%%%%%%%%%%%%%%%%%%%%%%%%%%%%%%%%%%%%%
\begin{frame}[fragile]\frametitle{तन्मात्र Tanmatras and Sense Organs}
    \textbf{तन्मात्र Tanmatras: Fundamental/seed Elements}
	
	
    \begin{tabular}{|l|l|l|}
    \hline
    \textbf{Element} & \textbf{Sanskrit Word} & \textbf{English Translation} \\
    \hline
    \textbf{Space/Ether} & शब्द (Shabda) & Sound \\
    \hline
    \textbf{Air} & स्पर्श (Sparsha) & Touch \\
    \hline
    \textbf{Fire} & रूप (Rupa) & Shape/Visual/Light \\
    \hline
    \textbf{Water} & रस (Rasa) & Taste \\
    \hline
    \textbf{Earth} & गन्ध (Gandha) & Smell \\
    \hline
    \end{tabular}


    \textbf{पञ्च  ज्ञानेन्द्रिय  Pancha Jnanedriya: Five Sense Organs}
	
    \begin{tabular}{|l|l|}
    \hline
    \textbf{Element} & \textbf{Sense Organ} \\
    \hline
    \textbf{Space/Ether} & Ears \\
    \hline
    \textbf{Air} & Skin \\
    \hline
    \textbf{Fire} & Eyes \\
    \hline
    \textbf{Water} & Tongue \\
    \hline
    \textbf{Earth} & Nose \\
    \hline
    \end{tabular}
\end{frame}

%%%%%%%%%%%%%%%%%%%%%%%%%%%%%%%%%%%%%%%%%%%%%%%%%%%%%%%%%%%
\begin{frame}[fragile]\frametitle{पञ्च कर्मेन्द्रिय Pancha Karmendriya: Five Action Organs}
    \begin{tabular}{|l|l|}
    \hline
    \textbf{Action Organ} & \textbf{Function} \\
    \hline
    Tongue & Speech \\
    \hline
    Hand & Receiving and Holding \\
    \hline
    Feet & Movement \\
    \hline
    Genitals & Reproduction \\
    \hline
    Anus & Discharging the waste \\
    \hline
    \end{tabular}
\end{frame}

%%%%%%%%%%%%%%%%%%%%%%%%%%%%%%%%%%%%%%%%%%%%%%%%%%%%%%%%%%%
\begin{frame}[fragile]\frametitle{सप्त धातु  Sapta Dhatu: Seven Tissues}
    \begin{itemize}
        \item \textbf{रस  Rasa} - Plasma
        \item \textbf{रक्त  Rakta} - Blood
        \item \textbf{मांस  Mamsa} - Muscle
        \item \textbf{मेद  Meda} - Adipose
        \item \textbf{अस्थि Asthi} - Bone
        \item \textbf{मज्जा Majja} - Bone Marrow
        \item \textbf{शुक्र Shukra} - Reproductive Tissues
    \end{itemize}
    \vspace{0.5cm}
    \textbf{Function:} 
    
    The dhatus (धातु) build and hold up the body; each one feeds the next in turn (standard Ayurveda). The Guidebook says the food sheath is built of these seven.
    
\end{frame}

%%%%%%%%%%%%%%%%%%%%%%%%%%%%%%%%%%%%%%%%%%%%%%%%%%%%%%%%%%%
\begin{frame}[fragile]\frametitle{त्रिदोष  Tri Doshas: Three Doshas}
    \begin{tabular}{|l|l|}
    \hline
    \textbf{दोष Dosha} & \textbf{Properties} \\
    \hline
    \textbf{वात  Vata} & Dry, Light, Cold, Rough, Subtle, Mobile \\
    \hline
    \textbf{पित्त Pitta} & Slightly oily, Sharp, Hot, Light, Foul-smelling, Flowing, Liquid \\
    \hline
    \textbf{कफ  Kapha} & Unctuous, Cold, Heavy, Slow, Smooth, Slimy, Steady \\
    \hline
    \end{tabular}
    \vspace{0.5cm}
    \textbf{दोष Dosha Meaning:}

    दोष (Dosha) means that which spoils or goes out of balance. {\tiny (Ref: Ashtanga Hridaya, Sutrasthana)}
    
\end{frame}

%%%%%%%%%%%%%%%%%%%%%%%%%%%%%%%%%%%%%%%%%%%%%%%%%%%%%%%%%%
\begin{frame}[fragile]\frametitle{पञ्च  प्राण  Pancha Prana: Five Pranas in the Body}

	\begin{itemize} 
	\item \textbf{प्राण (Prana)}: Heart region.
	\item \textbf{अपान (Apana)}: Anal region.
	\item \textbf{उदान (Udana)}: Throat region.
	\item \textbf{समान (Samana)}: Navel region.
	\item \textbf{व्यान (Vyana)}: The whole body.
	\end{itemize}
	{\tiny Locations vary between texts; these follow the Gheranda Samhita.}

	\vspace{0.5cm} 
	
	\textbf{उपप्राण (Upapranas):} 
	\begin{itemize} 
	\item \textbf{नाग (Naga)}: Belching.
	\item \textbf{कुर्म (Kurma)}: Blinking.
	\item \textbf{कृकल (Krikal)}: Hunger and thirst (sneezing in another translation).
	\item \textbf{देवदत्त (Devadatta)}: Yawning.
	\item \textbf{धनञ्जय (Dhananjaya)}: Stays in the body even after death.
	\end{itemize}
\end{frame}

%%%%%%%%%%%%%%%%%%%%%%%%%%%%%%%%%%%%%%%%%%%%%%%%%%%%%%%%%%
\begin{frame}[fragile]\frametitle{प्राण Pranas}
        \begin{center}
        % FIG ycb1-fig-pranas | now: images/ycb1_pranas.jpg | final: books/YCB1/images/new/ycb1_pranas_new.pdf | register: books/YCB1/docs/image_register.csv
        % PROMPT SP-DIAGRAM (common text: books/YCB1/docs/image_system_prompt.md) | Used in block ycb1-th-principles-of-yoga, frame 'Pranas'. Concept diagram: Pranas. Boxes, circles and arrows only, drawn from the structure given in the notes block (the names and their order are listed there); Sanskrit terms in Roman letters with the English meaning in brackets, nothing else added. Flat colours that still read in grayscale. Best redrawn later in TikZ. Aspect ratio 4:3. Check before use: author and the Sanskrit check in the notes.
        % NEW-PIC approved by the author (2026-11-27); the original is commented out below it
        \includegraphics[width=0.6\linewidth,keepaspectratio]{new/ycb1_pranas_new.pdf}
        % original, replaced (author approved 2026-11-27): \includegraphics[width=0.6\linewidth,keepaspectratio]{ycb1_pranas}
				
		% credit of the replaced original: {\tiny (Ref: Raja Yoga Rishikesh)}
		{\tiny (Line diagram drawn for this deck.)}
        \end{center}	

\end{frame}
%%%%%%%%%%%%%%%%%%%%%%%%%%%%%%%%%%%%%%%%%%%%%%%%%%%%%%%%%%
\begin{frame}[fragile]\frametitle{चक्र  Chakras}
        \begin{center}
        % FIG ycb1-fig-chakras | now: images/ycb1_chakras.jpg | final: books/YCB1/images/new/ycb1_chakras_new.pdf | register: books/YCB1/docs/image_register.csv
        % PROMPT SP-DIAGRAM (common text: books/YCB1/docs/image_system_prompt.md) | Used in block ycb1-th-principles-of-yoga, frame 'Chakra Mudras (Energy Center Gestures); Chakras'. Concept diagram: Chakra Mudras (Energy Center Gestures); Chakras. Boxes, circles and arrows only, drawn from the structure given in the notes block (the names and their order are listed there); Sanskrit terms in Roman letters with the English meaning in brackets, nothing else added. Flat colours that still read in grayscale. Best redrawn later in TikZ. Aspect ratio 4:3. Check before use: author and the Sanskrit check in the notes.
        % NEW-PIC approved by the author (2026-11-27); the original is commented out below it
        \includegraphics[width=0.7\linewidth,keepaspectratio]{new/ycb1_chakras_new.pdf}
        % original, replaced (author approved 2026-11-27): \includegraphics[width=0.7\linewidth,keepaspectratio]{ycb1_chakras}
				
		% credit of the replaced original: {\tiny (Ref: Raja Yoga Rishikesh)}
		{\tiny (Line diagram drawn for this deck.)}
        \end{center}	

\end{frame}

%%%%%%%%%%%%%%%%%%%%%%%%%%%%%%%%%%%%%%%%%%%%%%%%%%%%%%%%%%
\begin{frame}[fragile]\frametitle{धातु Dhatus}
        \begin{center}
        % FIG ycb1-fig-dhatus | now: images/ycb1_dhatus.jpg | final: books/YCB1/images/new/ycb1_dhatus_new.pdf | register: books/YCB1/docs/image_register.csv
        % PROMPT SP-DIAGRAM (common text: books/YCB1/docs/image_system_prompt.md) | Used in block ycb1-th-principles-of-yoga, frame 'Dhatus'. Concept diagram: Dhatus. Boxes, circles and arrows only, drawn from the structure given in the notes block (the names and their order are listed there); Sanskrit terms in Roman letters with the English meaning in brackets, nothing else added. Flat colours that still read in grayscale. Best redrawn later in TikZ. Aspect ratio 4:3. Check before use: author and the Sanskrit check in the notes.
        % NEW-PIC approved by the author (2026-11-27); the original is commented out below it
        \includegraphics[width=0.5\linewidth,keepaspectratio]{new/ycb1_dhatus_new.pdf}
        % original, replaced (author approved 2026-11-27): \includegraphics[width=0.5\linewidth,keepaspectratio]{ycb1_dhatus}
				
		% credit of the replaced original: {\tiny (Ref: Param Yoga)}
		{\tiny (Line diagram drawn for this deck.)}
        \end{center}	

\end{frame}


%%%%%%%%%%%%%%%%%%%%%%%%%%%%%%%%%%%%%%%%%%%%%%%%%%%%%%%%%%%%%%%%%%%%%%%%%%%%%%%%%
\begin{frame}[fragile]\frametitle{}
\begin{center}
{\Large 1.5 Introduction to major schools of Yoga (ज्ञान Jnana, भक्ति  Bhakti, कर्म  Karma, पतञ्जलि  Patanjali, हठ Hatha)}
\end{center}
\end{frame}

%%%%%%%%%%%%%%%%%%%%%%%%%%%%%%%%%%%%%%%%%%%%%%%%%%%%%%%%%%
\begin{frame}[fragile]\frametitle{Philosophic Schools in Ancient India}
    \textbf{Types of Schools (दर्शन  Darshana):}
    \begin{itemize}
        \item \textbf{आस्तिक  दर्शन Astik Darshana} - Schools that follow the Vedas. (वेद प्रमाण ): न्याय (गौतम), वैशेषिक (कणाद), सांख्य (कपिल), मीमांसा (जैमिनी), योग(पतंजली), वेदांत (बादरायण)
        \item \textbf{नास्तिक  दर्शन  Nastik Darshana} - Schools that do not follow the authority of the Vedas नास्तिक : जैन (महावीर), बौद्ध , चार्वाक (बृहस्पती)
    \end{itemize}
	
    \textbf{Paths of Yoga:}
    \begin{itemize}
        \item \textbf{ज्ञान योग Jnana Yoga}
        \item \textbf{भक्ती योग  Bhakti Yoga}
        \item \textbf{कर्म योग Karma Yoga}
        \item \textbf{राज योग Raja Yoga}
        \item \textbf{हठ योग Hatha Yoga}
        \item \textbf{मंत्र योग Mantra Yoga}
    \end{itemize}	
\end{frame}

% %%%%%%%%%%%%%%%%%%%%%%%%%%%%%%%%%%%%%%%%%%%%%%%%%%%%%%%%%%
% \begin{frame}[fragile]\frametitle{Astik Darshanas: Schools Following the Vedas}
    % \begin{tabular}{|l|l|}
    % \hline
    % \textbf{School Name} & \textbf{Founder} \\
    % \hline
    % \textbf{Nyaya} & Gautama Rishi \\
    % \hline
    % \textbf{Vaiseshika} & Kanada Rishi \\
    % \hline
    % \textbf{Sankhya} & Kapila Muni \\
    % \hline
    % \textbf{Yoga} & Patanjali Maharishi \\
    % \hline
    % \textbf{Purva Mimamsa} & Jaimini Rishi \\
    % \hline
    % \textbf{Uttara Mimamsa or Vedanta} & Vyasa Maharishi \\
    % \hline
    % \end{tabular}
% \end{frame}

% %%%%%%%%%%%%%%%%%%%%%%%%%%%%%%%%%%%%%%%%%%%%%%%%%%%%%%%%%%
% \begin{frame}[fragile]\frametitle{Nastik Darshanas: Schools Not Following the Vedas}
    % \begin{itemize}
        % \item \textbf{Charvaka}
        % \item \textbf{Buddhism}
        % \item \textbf{Jaina}
    % \end{itemize}
% \end{frame}

%%%%%%%%%%%%%%%%%%%%%%%%%%%%%%%%%%%%%%%%%%%%%%%%%%%%%%%%%%
\begin{frame}[fragile]\frametitle{ज्ञान योग  Jnana Yoga: Yoga of Wisdom}
    \textbf{Definition:} ज्ञान योग  Jnana Yoga is the path of knowledge: telling the real from the unreal. It is strictly non-dual (Advaita Vedanta) and contemplative.


    \textbf{Three Stages of ज्ञान योग  Jnana Yoga Practice:}
    \begin{itemize}
        \item \textbf{श्रवण Sravana} - Listening to the teaching
        \item \textbf{मनन Manana} - Reflecting and reasoning on it
        \item \textbf{निदिध्यासन  Nididhyasana} - Meditating on it again and again until the conviction is firm
    \end{itemize}

    \textbf{साधना चतुस्तय Sadhana Chatustaya (fourfold preparation):}
    \begin{itemize}
        \item \textbf{विवेक Viveka} - Discrimination of real from unreal
        \item \textbf{वैराग्य Vairagya} - Dispassion
        \item \textbf{षड्संपत Shadsampat} - Sixfold virtues
        \item \textbf{मुमुक्षत्व Mumukshutva} - Yearning for liberation
    \end{itemize}
	
The sixfold virtues are शम  Shama, दम  Dama, उपरति  Uparati, तितिक्षा Titiksha, श्रद्धा Sraddha and समाधान Samadhana.
	
\end{frame}


%%%%%%%%%%%%%%%%%%%%%%%%%%%%%%%%%%%%%%%%%%%%%%%%%%%%%%%%%%%
\begin{frame}[fragile]\frametitle{भक्ति  योग Bhakti Yoga: Yoga of Devotion}
    \textbf{Definition:} Bhakti Yoga is the path of devotion: unconditional, intense love for God.


    \textbf{Key Text:} नारद भक्ति सूत्र (Narada Bhakti Sutra) - a leading text on भक्ति Bhakti and प्रेम Prema (divine love).

    
    \textbf{Nine forms of devotion of भक्ति  योग Bhakti Yoga (Bhagavata Purana 7.5.23):}
    \begin{itemize}
        \item श्रवणं (shravan), कीर्तन (kirtan; the Guidebook calls it perhaps the most important)
        \item स्मरण् (smaran), पादसेवनम् (paadsevan)
        \item अर्चनम् (archanam), वन्दनम् (vandanam)
        \item दास्य (dasya), सख्यं (sakhya)
        \item आत्म निवेदनम् (aatma nivedan)
    \end{itemize}

    
    \textbf{Types of भक्त  Bhaktas (Devotees) in भगवद्गीता  Bhagavad Gita (7.16):}
    \begin{itemize}
        \item \textbf{आर्त  Arta} - Distressed
        \item \textbf{अर्थार्थी Artharthee} - Desirer of wealth
        \item \textbf{जिज्ञासु Jijnasu} - Inquisitive
        \item \textbf{ज्ञानी Jnani} - Knowledgeable
		
    \end{itemize}
\end{frame}

%%%%%%%%%%%%%%%%%%%%%%%%%%%%%%%%%%%%%%%%%%%%%%%%%%%%%%%%%%%
\begin{frame}[fragile]\frametitle{कर्म योग  Karma Yoga: Yoga of Action}
          \begin{itemize}
	  \item \textbf{Definition:} Karma Yoga is the path of action: work without attachment to the result, offering the fruit to God.


	  \item \textbf{Core Principles of Karma Yoga (According to भगवद्गीता Bhagavad Gita):}
    \begin{enumerate}
        \item Act from a sense of duty.
        \item Stay free of strong attachment to the outcome.
        \item Do not let worry about results disturb the mind while working.
        \item Meet success and failure with an even mind.
    \end{enumerate}

	  \item  \textbf{Objective:} To melt the ego into selfless action.
	      \end{itemize}

\end{frame}

%%%%%%%%%%%%%%%%%%%%%%%%%%%%%%%%%%%%%%%%%%%%%%%%%%%%%%%%%%%
\begin{frame}[fragile]\frametitle{राज  योग  Raja Yoga: The Royal Yoga}
    
      \begin{itemize}
	  \item \textbf{Definition:} Raja Yoga means the royal yoga; it is more the final state of Self-realisation than a separate technique.
	  \item  \textbf{Significance:} The हठयोग प्रदीपिका (Hatha Yoga Pradipika) calls hatha yoga a staircase leading to raja yoga.
	  \item \textbf{Objective:} To still the movements of the mind, so that the true self stands in its purity (Patanjali).
		\item Based on the eight limbs of Yoga (\textit{अष्टान्ग Ashtanga Yoga}).
		\item Major text: \textit{Yoga Sutras of Patanjali}.
		\item Raja yoga works on the mind; hatha yoga works on the body.
	  \end{itemize}
    
\end{frame}

%%%%%%%%%%%%%%%%%%%%%%%%%%%%%%%%%%%%%%%%%%%%%%%%%%%%%%%%%%%
\begin{frame}[fragile,shrink=6]\frametitle{हठ योग Hatha Yoga: The Physically Oriented Yoga}
    
\begin{itemize}
    \item Physically oriented; the Guidebook says it is sourced from Tantra.
    \item Most hatha systems begin with purification (\textit{Shodhana Kriyas}/\textit{Shat Kriyas}, \textit{शोधन क्रिया}/\textit{षट्क्रिया}); Patanjali's eight limbs begin with \textit{Yamas} (\textit{यम}) and \textit{Niyamas} (\textit{नियम}). This is about where practice starts, not a rejection of them.
    \item Core texts:
    \begin{itemize}
        \item \textit{Hatha Yoga Pradipika} (\textit{हठयोगप्रदीपिका}) by Swatmarama
        \item \textit{Gheranda Samhita} (\textit{घेरंड संहिता}) by Sage Gheranda
        \item \textit{Goraksha Samhita} (\textit{गोरक्ष संहिता}) by Gorakshanath
        \item \textit{Shiv Samhita}, \textit{Yoga Ratnavali}, \textit{Yoga Tarangini}, and others.
    \end{itemize}
    \item ``Hatha'' (\textit{हठ}): ``ha'' (\textit{ह}, sun, \textit{सूर्य}) and ``tha'' (\textit{ठ}, moon, \textit{चन्द्र}): a union of opposites.
    \item Sun: heat, masculinity, effort; Moon: coolness, femininity, surrender.
    \item Hatha yoga balances these opposites, aiding self-discovery.
    \item Misconception: ``hatha'' does not mean stubbornness; it stands for union.
    \item \textit{Hatha Pradipika} (\textit{हठयोगप्रदीपिका}): the sun channel \textit{Pingala Nadi} (\textit{पिंगला नाड़ी}) and moon channel \textit{Ida Nadi} (\textit{इडा नाड़ी}) unite, and \textit{Kundalini} (\textit{कुण्डलिनी}) awakens.
    \textit{Kundalini} rises through \textit{Sushumna Nadi} (\textit{सुषुम्ना नाड़ी}), crosses the \textit{Chakras} (\textit{चक्र}) and unites with \textit{Brahman} (\textit{ब्रह्मन्}) at \textit{Sahasrara} (\textit{सहस्रार}).
    \item Union of \textit{Atma} (\textit{आत्मा}) and \textit{Paramatma} (\textit{परमात्मा}), leading to enlightenment.
    \item Goal: control body and mind as a base for spiritual practice.
\end{itemize}
    
\end{frame}


%%%%%%%%%%%%%%%%%%%%%%%%%%%%%%%%%%%%%%%%%%%%%%%%%%%%%%%%%%%
\begin{frame}[fragile]\frametitle{मंत्र योग Mantra Yoga: The Chanting oriented Yoga}
    
\begin{itemize}
    \item Mantra Yoga is the Yoga of sound.
    \item The word \textit{mantra मन्त्र} is explained as ``man'' (mind, मन) and ``tra'' (to protect, त्र).
    \item \textit{Mantra} is that which protects the one who reflects on it. मननात्रायते इति मन्त्रः
    \item A mantra is a sacred sound with spiritual power.
    \item Mantra Yoga uses such sound for deep meditation.
    % "the power of the mantra is realised through Mantra Yoga": restates the line above, no separate point in the notes
    % \item The power of the \textit{mantra} is realized through Mantra Yoga.
    \item The best known mantra is \textbf{Om} (\textit{ॐ}).
    \item \textbf{OM} is called the \textit{mahamantra} (\textit{महामन्त्र}), the great mantra.
\end{itemize}
    
\end{frame}



%%%%%%%%%%%%%%%%%%%%%%%%%%%%%%%%%%%%%%%%%%%%%%%%%%%%%%%%%%%%%%%%%%%%%%%%%%%%%%%%%%
\begin{frame}[fragile]\frametitle{}
\begin{center}
{\Large 1.6 Introduction to Yoga practices for health and well being}
\end{center}
\end{frame}


%%%%%%%%%%%%%%%%%%%%%%%%%%%%%%%%%%%%%%%%%%%%%%%%%%%%%%%%%%%
\begin{frame}[fragile]\frametitle{Principles of Yoga Life}
योग्य जीवन के पांच पड़ाव 
      \begin{itemize}
		\item  व्यायाम: आसन 
		\item  श्वास: प्राणायाम 
		\item  आराम : शवासन 
		\item  अन्न : शाकाहारी 
		\item  सकारात्मक विचार, ध्यान 
	  \end{itemize}

\end{frame}
%%%%%%%%%%%%%%%%%%%%%%%%%%%%%%%%%%%%%%%%%%%%%%%%%%%%%%%%%%%
\begin{frame}[fragile]\frametitle{Health: Meaning and Definition}
    \textbf{World Health Organization (WHO) Definition:}
    
        Health is a state of complete physical, mental and social well-being and not merely the absence of disease or infirmity (WHO Constitution, 1948). The Guidebook adds spiritual well-being.
    

    
    \textbf{Sanskrit Definition:} 
    
        Health is स्वास्थ्य (Swasthya): स्व (Swa, self) and स्थ (Stha, abiding), so abiding in one's own Self. The Guidebook describes that Self as सत्-चित्-आनन्द (Sat-Chit-Ananda): being, consciousness, bliss.
    

    
    \textbf{Health according to Yoga (Guidebook):}
    \begin{itemize}
        \item Relaxed muscles
        \item Loose joints that save energy
        \item Low metabolic rate at rest
        \item Efficient use of energy
        \item Organ systems that stay coordinated even under stress
    \end{itemize}
\end{frame}

%%%%%%%%%%%%%%%%%%%%%%%%%%%%%%%%%%%%%%%%%%%%%%%%%%%%%%%%%%%
\begin{frame}[fragile]\frametitle{Strength and Balance in Yoga}
    \textbf{Strength:}
    \begin{itemize}
        \item नौकासन (Naukasana), उत्कटासन (Utkatasana) and भुजंगासन (Bhujangasana) build muscle strength, much as push-ups, lunges or squats do.
    \end{itemize}


    \textbf{Balance:}
    \begin{itemize}
        % not in the notes (ycb1-th-practices-for-health); kept as a teaching remark
        \item Balance is part of fitness and is easy to overlook in a gym routine.
        \item वृक्षासन (Vrikshasana, tree pose) trains standing steady on one leg, and so builds balance.
    \end{itemize}
\end{frame}

%%%%%%%%%%%%%%%%%%%%%%%%%%%%%%%%%%%%%%%%%%%%%%%%%%%%%%%%%%%
\begin{frame}[fragile]\frametitle{Joint Mobility and Yogic View of Health}
    \textbf{Joint Mobility:}
    \begin{itemize}
        \item Many poses move the joints in ways that walking or running do not.
        \item This widens the range of movement.
        % removed: "helps prevent injuries, reducing stiffness and pain" is a health claim the notes do not support
    \end{itemize}


    \textbf{Yogic View of Health:}
    \begin{itemize}
        \item श्वेताश्वतर उपनिषद् (Shvetasvatara Upanishad): first signs of progress are a light body, health, freedom from craving, clear skin, a pleasant voice and odour, and scanty excretions.
        \item Patanjali's योग दर्शन (Yoga Darshan, Vibhuti Pada): a perfected body has beauty, grace, strength and adamantine hardness; mastery of सामान वायु (Samana Vayu) brings radiance (3.40 in standard editions).
    \end{itemize}
\end{frame}

% %%%%%%%%%%%%%%%%%%%%%%%%%%%%%%%%%%%%%%%%%%%%%%%%%%%%%%%%%%%
% \begin{frame}[fragile]\frametitle{Mental and Emotional Well-being}
    % \textbf{Mental Well-being:}
    % \begin{itemize}
        % \item Conscious process of calming the mind removes weaknesses.
        % \item Builds willpower.
        % \item Steadfastness becomes the core.
        % \item Converts challenges into opportunities.
    % \end{itemize}

    
    % \textbf{Emotional Well-being:}
    % \begin{itemize}
        % \item Overcomes emotional conflicts, complexes, and turmoils.
        % \item Sharpens and sensitizes emotions systematically.
    % \end{itemize}
% \end{frame}


%%%%%%%%%%%%%%%%%%%%%%%%%%%%%%%%%%%%%%%%%%%%%%%%%%%%%%%%%%% 
\begin{frame}[fragile]\frametitle{Yoga Practices for Holistic Well-Being}

	\begin{itemize} 
	\item \textbf{Physical Health} 
		\begin{itemize} 
			\item \textit{Asanas} (आसन) - Steady and strengthen the body.
			\item \textit{Pranayama} (प्राणायाम) - Traditionally said to calm and lighten.
			\item \textit{Shavasana} (शवासन) - Rest; relieves tension.
			\item \textit{Kriyas} (क्रियाः) - Traditionally used for cleansing the body.
			\item \textit{Bandhas} (बन्धाः) (Body Locks) - Linked with pranayama; traditionally said to help control the mind.
		\end{itemize}

	\item \textbf{Mental Well-Being}
		\begin{itemize}
			\item \textit{Meditation} (Dhyana) - Steadies the mind.
			\item \textit{Prayer} - A supportive practice for mental well-being.
			\item \textit{Mantra Chanting} - Traditionally used to calm the mind.
		\end{itemize}

	\item \textbf{Emotional Balance}
		\begin{itemize}
			\item \textit{Bhakti Yoga} - Cultivates emotional connection and devotion.
		\end{itemize}

	\item \textbf{Overall Well-Being}
		\begin{itemize}
			\item \textit{Balanced Routine} - Make yoga part of daily life.
			\item \textit{Personalized Practice} - The teacher picks the mix for each person.
		\end{itemize}
	\end{itemize}

\end{frame}


% %%%%%%%%%%%%%%%%%%%%%%%%%%%%%%%%%%%%%%%%%%%%%%%%%%%%%%%%%%% 
% \begin{frame}[fragile]\frametitle{Yoga Practices for Holistic Well-Being}

	% \begin{itemize} 
	% \item \textbf{Physical Health:} 
		% \begin{itemize} 
			% \item \textit{Asanas} (आसन) - Enhance flexibility, strength, and balance. 
			% \item \textit{Pranayama} (प्राणायाम) - Improves lung capacity and energy flow. 
			% \item \textit{Shavasana} (शवासन) - Promotes relaxation and recovery. 
			% \item \textit{Kriyas} (क्रियाः) - Detoxify and purify the body. 
			% \item \textit{Bandhas} (बन्धाः) (Body Locks) - Support internal organ function and stability.
		% \end{itemize}

	% \item \textbf{Mental Well-Being:}
		% \begin{itemize}
			% \item \textit{Meditation} (Dhyana) - Reduces stress and enhances focus.
			% \item \textit{Mindfulness} - Increases present-moment awareness.
			% \item \textit{Mantra Chanting} - Provides mental clarity.
		% \end{itemize}

	% \item \textbf{Emotional Balance:}
		% \begin{itemize}
			% \item \textit{Bhakti Yoga} - Cultivates emotional connection and devotion.
			% \item \textit{Gratitude Practice} - Enhances positive emotional states.
		% \end{itemize}

	% \item \textbf{Overall Well-Being:}
		% \begin{itemize}
			% \item \textit{Balanced Routine} - Integrate yoga into daily life.
			% \item \textit{Personalized Practice} - Adapt practices to individual needs.
		% \end{itemize}
	% \end{itemize}

% \end{frame}	
			
% %%%%%%%%%%%%%%%%%%%%%%%%%%%%%%%%%%%%%%%%%%%%%%%%%%%%%%%%%%%
% \begin{frame}[fragile]\frametitle{Yoga Practices for Physical Health}

      % \begin{itemize}
        % \item \textit{Asanas} (आसन) (Postures) - Enhance flexibility, strength, and balance.
        % \item \textit{Pranayama} (प्राणायाम) (Breath Control) - Improves lung capacity and energy flow.
        % \item \textit{Shavasana} (शवासन) (Corpse Pose) - Promotes relaxation and recovery.
        % \item \textit{Kriyas} (क्रियाः) (Cleansing Techniques) - Detoxify and purify the body.
        % \item \textit{Bandhas} (बन्धाः) (Body Locks) - Support internal organ function and stability.
      % \end{itemize}

% \end{frame}

% %%%%%%%%%%%%%%%%%%%%%%%%%%%%%%%%%%%%%%%%%%%%%%%%%%%%%%%%%%%
% \begin{frame}[fragile]\frametitle{Yoga Practices for Mental Well-Being}

      % \begin{itemize}
		% \item \textit{Meditation} (Dhyana) - Reduces stress and enhances focus.
		% \item \textit{Mindfulness} - Increases present-moment awareness and emotional stability.
		% \item \textit{Pranayama} - Balances the nervous system and calms the mind.
		% \item \textit{Mantra Chanting} - Provides mental clarity and emotional balance.
		% \item \textit{Yoga Nidra} - Deep relaxation practice that improves sleep quality.
	  % \end{itemize}

% \end{frame}

% %%%%%%%%%%%%%%%%%%%%%%%%%%%%%%%%%%%%%%%%%%%%%%%%%%%%%%%%%%%
% \begin{frame}[fragile]\frametitle{Yoga Practices for Emotional Balance}

      % \begin{itemize}
		% \item \textit{Bhakti Yoga} - Cultivates emotional connection and devotion.
		% \item \textit{Journaling} - Reflect on emotions and mental patterns.
		% \item \textit{Gratitude Practice} - Enhances positive emotional states.
		% \item \textit{Group Yoga Classes} - Builds community and support.
		% \item \textit{Breath Awareness} - Helps in managing emotional responses.
	  % \end{itemize}

% \end{frame}

% %%%%%%%%%%%%%%%%%%%%%%%%%%%%%%%%%%%%%%%%%%%%%%%%%%%%%%%%%%%
% \begin{frame}[fragile]\frametitle{Yoga Practices for Overall Well-Being}

      % \begin{itemize}
		% \item \textit{Balanced Routine} - Integrate yoga into daily life for consistent benefits.
		% \item \textit{Healthy Lifestyle Choices} - Complement yoga with proper nutrition and hydration.
		% \item \textit{Holistic Approach} - Address physical, mental, and spiritual aspects.
		% \item \textit{Regular Practice} - Ensure regular engagement for sustained well-being.
		% \item \textit{Personalized Practice} - Adapt practices to individual needs and goals.
	  % \end{itemize}

% \end{frame}



%%%%%%%%%%%%%%%%%%%%%%%%%%%%%%%%%%%%%%%%%%%%%%%%%%%%%%%%%%%%%%%%%%%%%%%%%%%%%%%%%%
\begin{frame}[fragile]\frametitle{}
\begin{center}
{\Large 1.7 Introduction to \textit{Shatkarma} (षट्कर्म): meaning, purpose and their significance in \textit{Yoga Sadhana} (योग साधना)}
\end{center}
\end{frame}


% %%%%%%%%%%%%%%%%%%%%%%%%%%%%%%%%%%%%%%%%%%%%%%%%%%%%%%%%%%%
% \begin{frame}[fragile]\frametitle{Introduction to षटकर्म  Shatkarma}

      % \begin{itemize}
		% \item \textit{षटकर्म  Shatkarma} - Six purification techniques in Yoga.
		% \item Aimed at cleansing the body and mind for enhanced spiritual practice.
		% \item Prepares the practitioner for deeper practices like meditation and advanced postures.
		% \item Integrates physical, mental, and energetic purification.
		% \item Essential for holistic Yoga practice and overall health.
	  % \end{itemize}

% \end{frame}

% %%%%%%%%%%%%%%%%%%%%%%%%%%%%%%%%%%%%%%%%%%%%%%%%%%%%%%%%%%%
% \begin{frame}[fragile]\frametitle{Meaning and Purpose of षटकर्म Shatkarma}

      % \begin{itemize}
		% \item \textit{क्रिया Kriya} - Techniques to purify and balance the body.
		% \item \textit{Purpose} - Remove toxins, enhance vitality, and stabilize the mind.
		% \item \textit{Integration} - Facilitates deeper Yoga practices and spiritual growth.
		% \item \textit{Holistic Cleansing} - Addresses physical, mental, and energetic levels.
		% \item Essential for overcoming physical and mental obstructions.
	  % \end{itemize}

% \end{frame}

% %%%%%%%%%%%%%%%%%%%%%%%%%%%%%%%%%%%%%%%%%%%%%%%%%%%%%%%%%%%
% \begin{frame}[fragile]\frametitle{Significance of Shatkarma in Yoga Sadhana}

      % \begin{itemize}
		% \item \textit{Health Benefits} - Improves digestion, detoxifies, and boosts immunity.
		% \item \textit{Mental Clarity} - Reduces stress and mental clutter.
		% \item \textit{Energetic Balance} - Regulates the flow of vital energy (\textit{Prana}).
		% \item \textit{Spiritual Preparation} - Prepares the practitioner for higher states of consciousness.
		% \item \textit{Preventive Measures} - Aids in preventing diseases and imbalances.
	  % \end{itemize}

% \end{frame}

%%%%%%%%%%%%%%%%%%%%%%%%%%%%%%%%%%%%%%%%%%%%%%%%%%%%%%%%%%% 
\begin{frame}[fragile]\frametitle{Introduction to षटकर्म Shatkarma}

  \begin{itemize}
	\item \textit{षटकर्म Shatkarma} - Shat is six, karma is action: six cleansing practices.
	\item \textit{Purpose} - Clear waste from the internal organs that the body does not clear well on its own.
	\item \textit{Integration} - Done before pranayama and the higher practices.
	\item \textit{Balance} - Helps keep the three humours (kapha, pitta, vata) in balance; traditionally used for cleansing, not a cure.
	\item \textit{Caution} - The Hatha Pradipika gives them to those with excess fat or phlegm; for others pranayama is enough. Take medical advice first.
	\item \textit{Caution} - Never teach dhauti, basti or nauli to beginners; learn them from a qualified teacher.
	\item \textit{Spiritual Preparation} - A clean body lets the practitioner go further on the path.
	% removed: "addresses physical, mental and energetic levels" is not in the notes (ycb1-th-shatkarma-introduction)
	% \item \textit{Holistic Cleansing} - Addresses physical, mental, and energetic levels.
  \end{itemize}

\end{frame}

%%%%%%%%%%%%%%%%%%%%%%%%%%%%%%%%%%%%%%%%%%%%%%%%%%%%%%%%%%%
\begin{frame}[fragile]\frametitle{Overview of \textit{Shatkarma} (षट्कर्म) Techniques}

      \begin{itemize}
        \item \textit{Dhauti} (धौति) - The digestive tract up to the stomach (with water, a tube or a cloth).
        \item \textit{Basti} (बस्ति) - The colon, with water (jala) or air (sthala).
        \item \textit{Neti} (नेति) - The nasal passage: salt water (jala neti) or a thread (sutra neti).
        \item \textit{Trataka} (त्राटक) - The eyes: steady gazing at one point without blinking.
        \item \textit{Nauli} (नौलि) - The abdomen, by churning the abdominal muscles.
        \item \textit{Kapalabhati} (कपालभाती) - ``Skull shining'': quick, forceful out-breaths that clear the front of the head and the lungs.
      \end{itemize}
      {\small Order of the six as in the text; the order differs between schools. After neti, do kapalabhati to clear left-over water.}

\end{frame}



%%%%%%%%%%%%%%%%%%%%%%%%%%%%%%%%%%%%%%%%%%%%%%%%%%%%%%%%%%%%%%%%%%%%%%%%%%%%%%%%%%
\begin{frame}[fragile]\frametitle{}
\begin{center}
{\Large 1.8 Introduction to Yogic \textit{Sukshma Vyayama} (सूक्ष्म व्यायाम), \textit{Sthula Vyayama} (स्थूल व्यायाम) and \textit{Surya Namaskara} (सूर्य नमस्कार)}
\end{center}
\end{frame}

%%%%%%%%%%%%%%%%%%%%%%%%%%%%%%%%%%%%%%%%%%%%%%%%%%%%%%%%%%% 
\begin{frame}[fragile]\frametitle{Introduction to सूक्ष्म Sukshma and स्थूल Sthula व्यायामVyayama}

  \begin{itemize}
	\item \textit{सूक्ष्म Sukshma व्यायामVyayama} - Small joint and breathing movements, worked from head to toe.
	\item \textit{Enhances Mobility} - Keeps the joints supple; loosens, warms and relaxes the body.
	\item \textit{Prepares the Body} - Needs no special preparation and little time.
	\item \textit{स्थूल Sthula व्यायामVyayama} - Larger whole-body movements; the syllabus names sarvanga pushti and hrid gati (engine daud).
	\item \textit{General Fitness} - Like other regular activity, helps keep a healthy body weight (WHO).
	\item \textit{Comprehensive Practice} - Subtle and gross work together give balanced fitness.
  \end{itemize}
\end{frame}


% %%%%%%%%%%%%%%%%%%%%%%%%%%%%%%%%%%%%%%%%%%%%%%%%%%%%%%%%%%%
% \begin{frame}[fragile]\frametitle{Introduction to Sukshma Vyayama}

      % \begin{itemize}
		% \item \textit{Sukshma Vyayama} - Subtle exercises for the body and mind.
		% \item Focuses on gentle movements to prepare the body for more intensive practices.
		% \item Enhances joint mobility and flexibility.
		% \item Aids in the smooth functioning of internal organs.
		% \item Ideal for warming up and increasing energy flow.
	  % \end{itemize}

% \end{frame}

% %%%%%%%%%%%%%%%%%%%%%%%%%%%%%%%%%%%%%%%%%%%%%%%%%%%%%%%%%%%
% \begin{frame}[fragile]\frametitle{Purpose and Benefits of Sukshma Vyayama}

      % \begin{itemize}
		% \item \textit{Improves Circulation} - Enhances blood flow to muscles and joints.
		% \item \textit{Increases Flexibility} - Promotes flexibility in joints and muscles.
		% \item \textit{Prepares Body} - Warms up the body before more rigorous exercises.
		% \item \textit{Reduces Stiffness} - Alleviates joint and muscle stiffness.
		% \item \textit{Calms Mind} - Prepares the mind for focused practice.
	  % \end{itemize}

% \end{frame}

% %%%%%%%%%%%%%%%%%%%%%%%%%%%%%%%%%%%%%%%%%%%%%%%%%%%%%%%%%%%
% \begin{frame}[fragile]\frametitle{Introduction to Sthula Vyayama}

      % \begin{itemize}
		% \item \textit{Sthula Vyayama} - Gross or physical exercises for the body.
		% \item Includes more intense physical postures and movements.
		% \item Aims to build strength, endurance, and overall physical fitness.
		% \item Often used in combination with \textit{Sukshma Vyayama} for comprehensive practice.
		% \item Focuses on major muscle groups and physical conditioning.
	  % \end{itemize}

% \end{frame}

% %%%%%%%%%%%%%%%%%%%%%%%%%%%%%%%%%%%%%%%%%%%%%%%%%%%%%%%%%%%
% \begin{frame}[fragile]\frametitle{Purpose and Benefits of Sthula Vyayama}

      % \begin{itemize}
		% \item \textit{Strength Building} - Develops muscle strength and endurance.
		% \item \textit{Improves Posture} - Enhances overall body alignment and posture.
		% \item \textit{Boosts Fitness} - Increases cardiovascular and physical fitness.
		% \item \textit{Enhances Vitality} - Promotes general physical health and energy.
		% \item \textit{Supports Weight Management} - Aids in maintaining a healthy weight.
	  % \end{itemize}

% \end{frame}

%%%%%%%%%%%%%%%%%%%%%%%%%%%%%%%%%%%%%%%%%%%%%%%%%%%%%%%%%%% 
\begin{frame}[fragile]\frametitle{Introduction to \textit{Surya Namaskara} (सूर्य नमस्कार)}

  \begin{itemize}
	\item \textit{Surya Namaskara} (सूर्य नमस्कार) - Sun salutation: a flowing sequence of twelve positions with the breath, not a single asana.
	\item \textit{Purpose} - A salute to the sun and its energy.
	\item \textit{Integrates} - Movement, breath and intention.
	\item \textit{Loosens the Body} - Bends the spine forward and back; keeps spine and limbs supple.
	\item \textit{Supports Circulation} - And correct breathing.
	\item \textit{Prepares the Body} - Warms the body for asana and pranayama; traditionally credited with less belly fat.
	\item \textit{Timing} - By custom early morning or evening, facing the sun, on an empty stomach; custom, not a rule.
  \end{itemize}
\end{frame}

% %%%%%%%%%%%%%%%%%%%%%%%%%%%%%%%%%%%%%%%%%%%%%%%%%%%%%%%%%%%
% \begin{frame}[fragile]\frametitle{Introduction to \textit{Surya Namaskara} (सूर्य नमस्कार)}

      % \begin{itemize}
        % \item \textit{Surya Namaskara} (सूर्य नमस्कार) - Sun Salutation, a series of dynamic postures.
        % \item Traditionally performed to honor the Sun and its energy.
        % \item Consists of a sequence of 12 postures.
        % \item Integrates movement, breath, and intention.
        % \item Enhances overall physical and mental health.
      % \end{itemize}

% \end{frame}


% %%%%%%%%%%%%%%%%%%%%%%%%%%%%%%%%%%%%%%%%%%%%%%%%%%%%%%%%%%%
% \begin{frame}[fragile]\frametitle{Purpose and Benefits of Surya Namaskara}

      % \begin{itemize}
		% \item \textit{Improves Flexibility} - Stretches and tones the muscles.
		% \item \textit{Boosts Circulation} - Enhances blood flow and energy levels.
		% \item \textit{Increases Strength} - Builds core strength and endurance.
		% \item \textit{Balances Mind} - Calms the mind and prepares for meditation.
		% \item \textit{Energizes Body} - Invigorates and revitalizes overall health.
	  % \end{itemize}

% \end{frame}


%%%%%%%%%%%%%%%%%%%%%%%%%%%%%%%%%%%%%%%%%%%%%%%%%%%%%%%%%%%%%%%%%%%%%%%%%%%%%%%%%%
\begin{frame}[fragile]\frametitle{}
\begin{center}
{\Large 1.9 Introduction to \textit{Yogasana} (योगासन): meaning, principles, and their health benefits}
\end{center}
\end{frame}

%%%%%%%%%%%%%%%%%%%%%%%%%%%%%%%%%%%%%%%%%%%%%%%%%%%%%%%%%%% 
\begin{frame}[fragile]\frametitle{Introduction to Yogasana}

  \begin{itemize}
	\item \textit{Yogasana} - A way of sitting or holding the body; the Yoga Sutra asks for a steady, comfortable posture (sthira sukham asanam).
	\item \textit{Aim} - Poise, balance and strength; a calm mind ready for meditation.
	\item \textit{Principles} - Move slowly with the breath; hold steadily with awareness; no strain or comparing; end by relaxing.
	\item \textit{Health Benefits} - Calmness, flexibility, range of motion, strength and balance; for prevention and well-being, not a cure.
	\item \textit{Hatha Pradipika} - Health (arogya), steadiness (sthairya) and lightness of body (angalaghava).
	\item \textit{Calm Mind} - The Guidebook calls the physical gains incidental to mental calm.
  \end{itemize}

\end{frame}

% %%%%%%%%%%%%%%%%%%%%%%%%%%%%%%%%%%%%%%%%%%%%%%%%%%%%%%%%%%%
% \begin{frame}[fragile]\frametitle{Introduction to Yogasana}

      % \begin{itemize}
		% \item \textit{Yogasana} - Physical postures or poses in Yoga.
		% \item Derived from the Sanskrit words \textit{Yoga} (union) and \textit{Asana} (seat or posture).
		% \item Aims to harmonize body and mind through physical practice.
		% \item Forms the foundation for many Yoga practices and techniques.
		% \item Enhances physical and mental well-being.
	  % \end{itemize}

% \end{frame}

% %%%%%%%%%%%%%%%%%%%%%%%%%%%%%%%%%%%%%%%%%%%%%%%%%%%%%%%%%%%
% \begin{frame}[fragile]\frametitle{Principles of Yogasana}

      % \begin{itemize}
		% \item \textit{Alignment} - Proper positioning of body parts for effectiveness and safety.
		% \item \textit{Breath Awareness} - Coordinating breath with movement to enhance practice.
		% \item \textit{Balance} - Achieving physical and mental equilibrium in poses.
		% \item \textit{Stability} - Maintaining a steady and comfortable posture.
		% \item \textit{Mindfulness} - Being present and focused during practice.
	  % \end{itemize}

% \end{frame}

% %%%%%%%%%%%%%%%%%%%%%%%%%%%%%%%%%%%%%%%%%%%%%%%%%%%%%%%%%%%
% \begin{frame}[fragile]\frametitle{Health Benefits of Yogasana}

      % \begin{itemize}
		% \item \textit{Improves Flexibility} - Enhances range of motion in joints and muscles.
		% \item \textit{Builds Strength} - Develops muscle strength and endurance.
		% \item \textit{Enhances Posture} - Promotes proper alignment and balance.
		% \item \textit{Boosts Circulation} - Improves blood flow and cardiovascular health.
		% \item \textit{Reduces Stress} - Calms the mind and reduces anxiety levels.
		% \item \textit{Improves Digestion} - Stimulates digestive organs and enhances metabolism.
		% \item \textit{Enhances Mental Clarity} - Promotes focus, concentration, and mental well-being.
	  % \end{itemize}

% \end{frame}



%%%%%%%%%%%%%%%%%%%%%%%%%%%%%%%%%%%%%%%%%%%%%%%%%%%%%%%%%%%%%%%%%%%%%%%%%%%%%%%%%%
\begin{frame}[fragile]\frametitle{}
\begin{center}
{\Large 1.10 Introduction to \textit{Pranayama} (प्राणायाम) and \textit{Dhyana} (ध्यान) and their health benefits}
\end{center}
\end{frame}

%%%%%%%%%%%%%%%%%%%%%%%%%%%%%%%%%%%%%%%%%%%%%%%%%%%%%%%%%%% 
\begin{frame}[fragile]\frametitle{Introduction to Pranayama}

  \begin{itemize}
	\item \textit{Pranayama} - Prana is life force, ayama is extension or control: control of prana through the breath.
	\item \textit{Three parts} - Puraka (inhalation), kumbhaka (retention), rechaka (exhalation).
	\item \textit{Integrates} - Can be joined with postures and meditation.
	\item \textit{Lung Capacity} - Traditionally credited with better breathing capacity; gains are listed for each practice.
	\item \textit{Calms the Mind} - Traditionally said to ease anxiety and stress.
	\item \textit{Relaxation} - Prepares the mind for meditation and concentration; not a treatment.
	\item \textit{Safety} - Breath holding and forceful breathing are not for beginners.
  \end{itemize}

\end{frame}

% %%%%%%%%%%%%%%%%%%%%%%%%%%%%%%%%%%%%%%%%%%%%%%%%%%%%%%%%%%%
% \begin{frame}[fragile]\frametitle{Introduction to Pranayama}

      % \begin{itemize}
		% \item \textit{Pranayama} - The practice of breath control in Yoga.
		% \item Derived from Sanskrit words \textit{Prana} (life force or breath) and \textit{Ayama} (control or extension).
		% \item Aims to regulate and expand the breath to enhance life energy.
		% \item Integrates breath with physical postures and meditation.
		% \item Fundamental for balancing mind and body.
	  % \end{itemize}

% \end{frame}

% %%%%%%%%%%%%%%%%%%%%%%%%%%%%%%%%%%%%%%%%%%%%%%%%%%%%%%%%%%%
% \begin{frame}[fragile]\frametitle{Health Benefits of Pranayama}

      % \begin{itemize}
		% \item \textit{Improves Lung Capacity} - Enhances respiratory efficiency and endurance.
		% \item \textit{Balances Nervous System} - Regulates stress and anxiety levels.
		% \item \textit{Enhances Concentration} - Increases mental focus and clarity.
		% \item \textit{Boosts Immunity} - Strengthens the immune system and overall vitality.
		% \item \textit{Aids in Detoxification} - Promotes the removal of toxins from the body.
		% \item \textit{Calms the Mind} - Reduces mental agitation and promotes relaxation.
		% \item \textit{Regulates Emotions} - Helps in managing emotional responses and stability.
	  % \end{itemize}

% \end{frame}

% %%%%%%%%%%%%%%%%%%%%%%%%%%%%%%%%%%%%%%%%%%%%%%%%%%%%%%%%%%%
% \begin{frame}[fragile]\frametitle{Introduction to Dhyana}

      % \begin{itemize}
		% \item \textit{Dhyana} - Meditation or the practice of focused attention.
		% \item Derived from Sanskrit meaning "profound contemplation" or "meditative absorption."
		% • Aims to achieve a state of mental stillness and clarity.
		% \item Involves sustained concentration and mindfulness.
		% \item Integral for achieving higher states of consciousness and inner peace.
	  % \end{itemize}

% \end{frame}

% %%%%%%%%%%%%%%%%%%%%%%%%%%%%%%%%%%%%%%%%%%%%%%%%%%%%%%%%%%%
% \begin{frame}[fragile]\frametitle{Health Benefits of Dhyana}

      % \begin{itemize}
		% \item \textit{Reduces Stress} - Lowers cortisol levels and promotes relaxation.
		% \item \textit{Enhances Emotional Well-Being} - Improves mood and emotional resilience.
		% \item \textit{Improves Focus} - Increases attention span and cognitive function.
		% \item \textit{Balances Blood Pressure} - Helps in maintaining healthy blood pressure levels.
		% \item \textit{Promotes Inner Peace} - Fosters a sense of calm and tranquility.
		% \item \textit{Aids in Self-Realization} - Encourages deeper self-awareness and understanding.
		% \item \textit{Supports Mental Health} - Helps in managing anxiety, depression, and other mental health issues.
	  % \end{itemize}

% \end{frame}


%%%%%%%%%%%%%%%%%%%%%%%%%%%%%%%%%%%%%%%%%%%%%%%%%%%%%%%%%%% 
\begin{frame}[fragile]\frametitle{Introduction to Dhyana}
  \begin{itemize}
	\item \textit{Dhyana} - Unbroken flow of attention toward one object, when concentration (dharana) continues without a break; it leads to samadhi.
	\item \textit{Purpose} - Calm and self-realisation (Guidebook).
	\item \textit{Well-being} - Traditionally credited with calm, clarity and concentration; a supportive practice, not a treatment.
  \end{itemize}
\end{frame}